% !TEX root = ../main.tex
\section{An accuracy-independent schedule with increasing \texorpdfstring{$s_k$}{s\_k}}
\label{sec:upper-bound}

This appendix gives a parameter schedule that does not depend on
\(\epsilon\). With \(\rho\) fixed, \(s_k=\Theta((k+1)^{1/3})\), and diminishing
center stepsizes, SPLALM finds an \(\epsilon\)-KKT point in
\(\mathcal O(\epsilon^{-3})\) iterations (Theorem~\ref{thm:main}). The
parameters vary with \(k\), so the schedule lies outside the family
\(\calC_\theta\) of Definition~\ref{def:algorithm-family}.

Throughout this section, we additionally impose the following assumption.
It is not part of the problem class \(\mathcal I(\mathbf C,\Delta)\) of
Definition~\ref{def:problem-class}.

\begin{assumption}[Level-bounded objective]
\label{ass:level-boundedness}
For every \(c\in\RR\), the sublevel set
\(\{\bx\in\RR^n:F(\bx)\le c\}\) is empty or compact.
\end{assumption}

\subsection{Algorithm with increasing \texorpdfstring{$s_k$}{s\_k}}
\label{sec:method}

This section uses the updates of Definition~\ref{def:splalm} with
iteration-dependent parameters \(s_k\), \(\eta_k\), \(\tau_k\), \(\alpha_k\),
and \(\beta_k\), and with an arbitrary initialization; see
Algorithm~\ref{alg:ddlalm}. By~\eqref{eq:reduced-al}--\eqref{eq:potential},
the first two updates read
\begin{equation}\label{eq:explicit-updates}
    \begin{aligned}
        \bx^{k+1}
        &=
        \bx^k-\eta_k
        \Bigl[
            \nabla f(\bx^k)
            +J_{\calA}(\bx^k)^\top\bq(\bx^k,\blmbd^k)
            +\frac{1}{\gamma}(\bx^k-\bz^k)
        \Bigr], \\
        \blmbd^{k+1}
        &=
        \blmbd^k+\tau_k
        \Bigl[
            \calA(\bx^{k+1})
            -\by(\bx^{k+1},\blmbd^k)
            -\frac{1}{s_k}(\blmbd^k-\bmu^k)
        \Bigr],
    \end{aligned}
\end{equation}
and the outputs \((\bx^{k+1},\by^{k+1},\bq^k)\) are those of
Definition~\ref{def:splalm}.

\begin{algorithm}[t]
    \caption{Smoothed-Perturbed Linearized Augmented Lagrangian Method (SPLALM)}
    \label{alg:ddlalm}
    \small
    \begin{algorithmic}[1]
        \STATE \textbf{Input:} Initial
        \(\bx^0,\bz^0\in\RR^n\) and
        \(\blmbd^0,\bmu^0\in\RR^m\); 
        Parameters: \(\rho>0\) and \(\gamma>0\);
        sequences: \(\{\eta_k\}\), \(\{\tau_k\}\), \(\{\alpha_k\}\), and
        \(\{\beta_k\}\);
        and iteration limit \(K\).
        \FOR{\(k=0,1,\ldots,K-1\)}
            \STATE  \(\bx^{k+1}=\bx^k- \eta_k \nabla_{\bx} \mathcal P_{s_k}(\bx^{k},\blmbd^{k};\bz^{k},\bmu^{k})\);
            \STATE  \(\blmbd^{k+1}=\blmbd^k+\tau_k
            \nabla_{\blmbd}\mathcal P_{s_k}
            (\bx^{k+1},\blmbd^{k};\bz^{k},\bmu^{k})\).
            \STATE 
            \(\bz^{k+1}=\bz^k+ \alpha_k(\bx^{k+1}-\bz^k)\).
            \STATE  \(\bmu^{k+1}=\bmu^k+ \beta_k(\blmbd^{k+1}-\bmu^k)\).
        \ENDFOR
    \end{algorithmic}
\end{algorithm}

We use fixed primal and dual stepsizes together with diminishing center
stepsizes. Choose \(\rho>0\) and \(\gamma>0\) satisfying
Proposition~\ref{prop:curvature-smoothness}. With \(L_x\) and
\(\mu_\gamma\) defined in Section~\ref{subsec:curvature-smoothness} and
\(L_{d,0}\) defined in Section~\ref{sec:complexity-analysis}, set \(\eta_k\equiv\eta\) and
\(\tau_k\equiv\tau\). An admissible choice is
\begin{equation}\label{eq:ddlalm-parameter-schedule}
\begin{aligned}
    s_k &= c_s(k+1)^{1/3},
    \qquad
    \eta = \frac{1}{2L_x},
    \qquad
    \tau =
    \left(
        2L_{d,0}
        +\frac{72L_xG_{\cal A}^2}{\mu_\gamma^2}
    \right)^{-1},
    \\
    \alpha_k &=
    \frac{\min\{s_0,\gamma\mu_\gamma\tau\}}{36s_k},
    \qquad
    \beta_k=\frac{\theta_\mu}{k+k_0},
\end{aligned}
\end{equation}
where \(c_s>0\), \(\theta_\mu\in(0,1]\), and \(k_0\ge1\) are constants
independent of \(K\). Consequently,
\[
    s_k=\Theta((k+1)^{1/3}),
    \qquad
    \alpha_k=\Theta((k+1)^{-1/3}),
    \qquad
    \beta_k=\Theta((k+1)^{-1}).
\]

\begin{theorem}
\label{thm:main}
Suppose that Assumptions~\ref{ass:level-boundedness},
\ref{ass:smoothness}, and~\ref{ass:h} hold.
Let the sequence generated by \Cref{alg:ddlalm} use
\eqref{eq:ddlalm-parameter-schedule}, with $\rho$ and $\gamma$ as specified
above. For $K\ge2$, choose a tail-best index
\[
    \kappa\in\argmin_{k\in\{\lceil K/2\rceil,\ldots,K-1\}}
    \left\{
        \|\Delta_{\bx}^{k+1}\|^2
        +\|\Delta_{\blmbd}^{k+1}\|^2
        +\alpha_k\|\bx^{k+1}-\bz^k\|^2
    \right\}.
\]
Then it holds that
\[
\max\left\{
\left[
    \left\|
        \nabla f({\bx}^{\kappa+1})
        +
        J_{\cal A}({\bx}^{\kappa+1})^\top{\bq}^{\kappa}
    \right\|^2
\right],
\;
\left[
    \left\|
        {\cal A}({\bx}^{\kappa+1})-{\by}^{\kappa+1}
    \right\|^2
\right]
\right\}
=
\mathcal O(K^{-2/3}),
\]
with \({\bq}^{\kappa}\in\partial h({\by}^{\kappa+1})\).
Consequently, given \(\epsilon>0\), the iteration complexity of
\Cref{alg:ddlalm} to obtain an \(\epsilon\)-KKT point
is of order \(\mathcal O(\epsilon^{-3})\).
\end{theorem}



\subsection{Complexity Analysis}
\label{sec:complexity-analysis}
\label{sec:basic-descent}
We use the constants \(L_{\calD,\bx}\), \(\mu_\gamma\), \(L_x\), and
\(L_\lambda(s)\) and Proposition~\ref{prop:curvature-smoothness} from
Section~\ref{subsec:curvature-smoothness}.
 
\subsubsection{Value Functions and Lyapunov Function}
\label{subsec:value-functions}

Fix $s>0$ and $(\bz,\bmu)$. By Proposition~\ref{prop:curvature-smoothness},
$\mathcal P_s(\cdot,\cdot;\bz,\bmu)$ is strongly convex in $\bx$ and strongly
concave in $\blmbd$. Hence the following value functions are well defined,
and the minimizer and maximizer below are unique:
\begin{align}
    d_s(\blmbd;\bz,\bmu)
    &:= \min_{\bx\in\RR^n}\mathcal P_s(\bx,\blmbd;\bz,\bmu),
    &
    \bx_s(\blmbd;\bz,\bmu)
    &:= \argmin_{\bx\in\RR^n}\mathcal P_s(\bx,\blmbd;\bz,\bmu),
    \label{eq:dual-value}\\
    p_s(\bz,\bmu)
    &:= \max_{\blmbd\in\RR^m} d_s(\blmbd;\bz,\bmu),
    &
    \blmbd_s^\star(\bz,\bmu)
    &:= \argmax_{\blmbd\in\RR^m} d_s(\blmbd;\bz,\bmu).
    \label{eq:saddle-value}
\end{align}
Strong convexity--strong concavity also yields strong duality,
\begin{equation}
    p_s(\bz,\bmu)
    = \min_{\bx\in\RR^n}\max_{\blmbd\in\RR^m}
      \mathcal P_s(\bx,\blmbd;\bz,\bmu),
    \label{eq:minimax-equality}
\end{equation}
and the unique saddle point is
$(\bx_s^\star(\bz,\bmu),\blmbd_s^\star(\bz,\bmu))$, where
\[
    \bx_s^\star(\bz,\bmu)
    :=\bx_s\bigl(\blmbd_s^\star(\bz,\bmu);\bz,\bmu\bigr).
\]

\blue{Following the three-value Lyapunov construction for smoothed nonconvex--concave minimax methods~\cite{zhang2020smoothed}, and its perturbed-smoothed extension~\cite{li2026smoothing}, we adapt the construction to the reduced augmented-Lagrangian setting and define}
\begin{equation}
\begin{aligned}
    \Psi_s(\bx,\blmbd;\bz,\bmu)
    &:={}
    \mathcal P_s(\bx,\blmbd;\bz,\bmu)
    -2d_s(\blmbd;\bz,\bmu)
    +2p_s(\bz,\bmu) \\
    &= \underbrace{\mathcal P_s(\bx,\blmbd;\bz,\bmu)-d_s(\blmbd;\bz,\bmu)}_{\text{primal gap at fixed }\blmbd}
     + \underbrace{p_s(\bz,\bmu)-d_s(\blmbd;\bz,\bmu)}_{\text{dual gap}}
     + p_s(\bz,\bmu).
\end{aligned}
\label{eq:lyapunov-decomposition}
\end{equation}
Both gaps are nonnegative, so $\Psi_s\ge p_s$.

\paragraph{Iteration-$k$ shorthand.}
At iteration $k$ we evaluate the above objects at $(s,\bmu)=(s_k,\bmu^k)$,
replace the subscript $s$ by $k$, and suppress $\bmu^k$. For example,
$\mathcal P_k(\bx,\blmbd;\bz):=\mathcal P_{s_k}(\bx,\blmbd;\bz,\bmu^k)$ and
$d_k(\blmbd;\bz):=d_{s_k}(\blmbd;\bz,\bmu^k)$; the quantities
$p_k$, $\bx_k$, $\blmbd_k^\star$, $\bx_k^\star$, and $\Psi_k$ are defined
analogously. We also write
$\Delta_{\bx}^{k+1}:=\bx^{k+1}-\bx^k$, and define
$\Delta_{\blmbd}^{k+1}$, $\Delta_{\bz}^{k+1}$, and $\Delta_{\bmu}^{k+1}$
in the same way.

\subsubsection{Basic One-Step Estimate}
\label{subsec:basic-one-step}

Throughout this subsection, we assume that 
\begin{equation}
    s_{k+1}\ge s_k>0,\qquad
    0 < \tau_k\le s_k,\qquad
    0<\beta_k\le1,\qquad
    0\le\alpha_k\le2.
    \label{eq:basic-parameter-condition}
\end{equation} 
The following lemma gives a basic one-step descent estimate of the potential function $\mathcal P$ from iteration $k$ to $k+1$.

\begin{lemma}
\label{lem:basic-one-step}
Suppose Assumptions~\ref{ass:smoothness} and~\ref{ass:h} hold.
Let the iterates be generated by Algorithm~\ref{alg:ddlalm} with the
parameters satisfying \eqref{eq:basic-parameter-condition}.  Then, for
every $k\ge0$, the following holds:
\begin{equation}
\label{eq:basic-one-step}
\begin{aligned}
    &\mathcal P_{k+1}(\bx^{k+1},\blmbd^{k+1};\bz^{k+1})
    -\mathcal P_k(\bx^k,\blmbd^k;\bz^k) \\
    \le\;&
    -\left(\frac{1}{\eta_k}-\frac{L_x}{2}\right)
      \|\Delta_{\bx}^{k+1}\|^2
    +\left(\frac{1}{\tau_k}-\frac{1}{2s_k}\right)
      \|\Delta_{\blmbd}^{k+1}\|^2 \\
    &-\frac{\alpha_k(2-\alpha_k)}{2\gamma}
      \|\bx^{k+1}-\bz^k\|^2
    +\frac{\beta_k(2-\beta_k)}{2s_k}
      \|\blmbd^{k+1}-\bmu^k\|^2 +\frac{1}{2}
      \left(\frac{1}{s_k}-\frac{1}{s_{k+1}}\right)
      \|\blmbd^{k+1}-\bmu^{k+1}\|^2 .
\end{aligned}
\end{equation}
\end{lemma}


As shown in Lemma~\ref{lem:basic-one-step}, the positive ascent term
generated by the $\blmbd$-update prevents the potential
$\mathcal P_k$ from yielding a direct descent estimate. We therefore relate
this term to the one-step increase of the dual value function, as developed
in Lemma~\ref{lem:dual-value-gain}.  This relation allows us to construct a Lyapunov function with a descent property. Before proceeding, we establish several regularity properties of the value functions \(d_k\) and \(p_k\), with their gradient formulas following from Danskin's theorem.
\begin{lemma}
\label{lem:v2-value-regularity}
Under the conditions of Proposition~\ref{prop:curvature-smoothness}, for all $k\ge0$, the following properties hold:
\begin{enumerate}[label=(\roman*)]
\item
For every \(\bz \in \RR^n \), the minimizer
\(\bx_k(\blmbd;\bz)\) is unique and
\(\frac{G_{\cal A}}{\mu_\gamma}\)-Lipschitz continuous in
\(\blmbd\).

\item for every \(\bz\in\RR^n\),
\(d_k(\cdot;\bz)\) is \(1/s_k\)-strongly concave and
\(L_{d,k}\)-smooth, with gradient
\[
    \nabla_{\blmbd}d_k(\blmbd;\bz)
    =
    \nabla_{\blmbd}\mathcal P_k
    \bigl(\bx_k(\blmbd;\bz),\blmbd;\bz\bigr),
\]
where \(L_{d,k}:=\frac1\rho+\frac1{s_k}+\frac{G_{\cal A}^2}{\mu_\gamma}\).

\item
The function \(p_k(z)\) is continuously differentiable, with gradient
\[
    \nabla p_k(\bz)
    =
    \frac1\gamma
    \bigl(\bz-\bx_k^\star(\bz)\bigr),
\]
and has a \(1/\gamma\)-Lipschitz continuous gradient.
\end{enumerate}
\end{lemma}

The following lemma controls the dual ascent term.
\begin{lemma}
\label{lem:dual-value-gain}
Suppose that the conditions of
Lemma~\ref{lem:v2-value-regularity} hold, and let the iterates
be generated by Algorithm~\ref{alg:ddlalm} with parameters satisfying
\eqref{eq:basic-parameter-condition}.
For every $k\ge 0$, assume that
\begin{equation}
\label{eq:dual-gain-condition}
    \frac{1}{\tau_k}+\frac{1}{2s_k} \ge 2L_{d,k}.
\end{equation}
Then the following holds
\begin{equation*}
\begin{aligned}
& \left(\frac{1}{\tau_k}-\frac{1}{2s_k}\right)
  \|\Delta_{\boldsymbol{\lambda}}^{k+1}\|^2 \\
& \le 2\!\bigl[
    d_k(\boldsymbol{\lambda}^{k+1};\mathbf{z}^k)
    -d_k(\boldsymbol{\lambda}^k;\mathbf{z}^k)
  \bigr]
  -\frac{L_{d,k}}{2}
    \|\Delta_{\boldsymbol{\lambda}}^{k+1}\|^2
  +\frac{4G_{\mathcal A}^2}{L_{d,k}\mu_\gamma}
  \left[
    \mathcal P_k(
      \mathbf{x}^{k+1},\boldsymbol{\lambda}^k;\mathbf{z}^k)
    -d_k(\boldsymbol{\lambda}^k;\mathbf{z}^k)
  \right].
\end{aligned}
\end{equation*}
Moreover, it holds that
\begin{equation}
\begin{aligned}
&\mathcal P_{k+1}(\bx^{k+1},\blmbd^{k+1};\bz^{k+1})
-\mathcal P_k(\bx^k,\blmbd^k;\bz^k)
\\
&\le
-\left[
\frac1{\eta_k}-\frac{L_x}{2}
-
\frac{2G_{\mathcal A}^2(1/\eta_k+L_x)^2}
{L_{d,k}\mu_\gamma^2}
\right]
\|\Delta_{\bx}^{k+1}\|^2
-\frac{L_{d,k}}2\|\Delta_{\blmbd}^{k+1}\|^2 \\
&\qquad
+2\bigl[d_k(\blmbd^{k+1};\bz^k)-d_k(\blmbd^k;\bz^k)\bigr]
-\frac{\alpha_k(2-\alpha_k)}{2\gamma}
\|\bx^{k+1}-\bz^k\|^2 \\
&\qquad
+\frac{\beta_k(2-\beta_k)}{2s_k}
\|\blmbd^{k+1}-\bmu^k\|^2
+\frac12\left(\frac1{s_k}-\frac1{s_{k+1}}\right)
\|\blmbd^{k+1}-\bmu^{k+1}\|^2.
\end{aligned}
\label{eq:dual-value-full-step}
\end{equation}
\end{lemma}

The potential $\mathcal P_k$ increases along the dual update, so its variation alone does not capture progress in
both variables. We therefore seek a merit function that accounts for both
primal suboptimality and progress in dual maximization. To account for
these two effects, consider the partial potential
\begin{equation}
\label{eq:partial-potential}
    \mathcal P_k-cd_k
    =(\mathcal P_k-d_k)-(c-1)d_k,
    \qquad c>1.
\end{equation}
The first term measures the primal optimality gap, while the second term
decreases as the dual value increases. In the following analysis, we set
$c=2$, consistent with the coefficient of the dual-value increment in
\eqref{eq:dual-value-full-step}.

\begin{lemma}
\label{lem:partial-lyapunov-z-step}
Under the conditions of Lemma~\ref{lem:dual-value-gain}, for every $k\ge0$, the following holds
\begin{equation}
\begin{aligned}
&\Big[
\mathcal P_{k+1}(\bx^{k+1},\blmbd^{k+1};\bz^{k+1})
-2d_k(\blmbd^{k+1};\bz^{k+1})
\Big] -
\Big[
\mathcal P_k(\bx^k,\blmbd^k;\bz^k)
-2d_k(\blmbd^k;\bz^k)
\Big]\\
\le\,&
-\left[
\frac1{\eta_k}-\frac{L_x}{2}
-\frac{2G_{\mathcal A}^2(1/\eta_k+L_x)^2}
{L_{d,k}\mu_\gamma^2}
\right]
\|\Delta_{\bx}^{k+1}\|^2 -\frac{L_{d,k}}{2}\|\Delta_{\blmbd}^{k+1}\|^2 \\
&\qquad
+\frac{\alpha_k(2-\alpha_k)}{2\gamma}
\|\bx^{k+1}-\bz^k\|^2
-\frac{2\alpha_k}{\gamma}
\left\langle
\bx^{k+1}
-\bx_k(\blmbd^{k+1};\bz^{k+1}),
\bx^{k+1}-\bz^k
\right\rangle  \\
&\qquad
+\frac{\beta_k(2-\beta_k)}{2s_k}
\|\blmbd^{k+1}-\bmu^k\|^2
+\frac12\left(\frac1{s_k}-\frac1{s_{k+1}}\right)
\|\blmbd^{k+1}-\bmu^{k+1}\|^2.
\end{aligned}
\label{eq:partial-lyapunov-z-step}
\end{equation}
\end{lemma}
Equation~\eqref{eq:partial-lyapunov-z-step} gives the one-step estimate
for the partial potential in \eqref{eq:partial-potential} with $c=2$.
The $\bx$- and
$\blmbd$-updates contribute the first two terms on the right-hand side.
The $\bz$-update, however, produces a positive quadratic term and an
indefinite inner-product term. Hence, the partial potential does not yet
provide a sufficient descent estimate for the full iteration.

At an exact saddle point of
\(\min_{\bx}\max_{\blmbd}\mathcal P_k(\bx,\blmbd;\bz)\),
strong duality gives
\[
    \mathcal P_k\bigl(\bx_k^\star(\bz),
    \blmbd_k^\star(\bz);\bz\bigr)
    =d_k\bigl(\blmbd_k^\star(\bz);\bz\bigr)
    =p_k(\bz),
\]
and therefore
\[
    \mathcal P_k\bigl(\bx_k^\star(\bz),
    \blmbd_k^\star(\bz);\bz\bigr)
    -2d_k\bigl(\blmbd_k^\star(\bz);\bz\bigr)
    =-p_k(\bz).
\]
Consequently, any $\bz$-update that decreases the saddle value $p_k$
increases the partial potential $\mathcal P_k-2d_k$. This motivates us
to examine whether an additional value-function term can restore a descent
estimate.

We therefore add the value-function increment to the partial potential in
\eqref{eq:partial-potential}:
$2[p_k(\bz^{k+1})-p_k(\bz^k)]$ and combine the two estimates directly.
The resulting one-step bound is stated next.

\begin{lemma}
\label{lem:z-step-p}
Under the conditions of Lemma~\ref{lem:v2-value-regularity} and
$0\le\alpha_k\le2$, the $\bz$-update satisfies
\begin{equation}
\begin{aligned}
p_k(\bz^{k+1})-p_k(\bz^k)
\le\,
-\frac{\alpha_k(2-\alpha_k)}{2\gamma}
\|\bx^{k+1}-\bz^k\|^2 +
\frac{\alpha_k}{\gamma}
\left\langle
\bx^{k+1}-\bx_k^\star(\bz^k),
\bx^{k+1}-\bz^k
\right\rangle.
\end{aligned}
\label{eq:z-step-p}
% \notag
\end{equation}
\end{lemma}

\begin{lemma}
\label{lem:augmented-one-step}
Under the conditions of Lemma~\ref{lem:partial-lyapunov-z-step}, for every
$k\ge0$,
\begin{equation}
\begin{aligned}
&\Big[
\mathcal P_{k+1}(\bx^{k+1},\blmbd^{k+1};\bz^{k+1})
-2d_k(\blmbd^{k+1};\bz^{k+1})
+2p_k(\bz^{k+1})
\Big] \\
&\quad-
\Big[
\mathcal P_k(\bx^k,\blmbd^k;\bz^k)
-2d_k(\blmbd^k;\bz^k)
+2p_k(\bz^k)
\Big] \\
&\le
-\left[
\frac1{\eta_k}-\frac{L_x}{2}
-\frac{2G_{\mathcal A}^2(1/\eta_k+L_x)^2}
{L_{d,k}\mu_\gamma^2}
\right]
\|\Delta_{\bx}^{k+1}\|^2
-\frac{L_{d,k}}{2}\|\Delta_{\blmbd}^{k+1}\|^2 \\
&\qquad
-\frac{\alpha_k(2-\alpha_k)}{2\gamma}
\|\bx^{k+1}-\bz^k\|^2
+\frac{2\alpha_k}{\gamma}
\left\langle
\bx_k(\blmbd^{k+1};\bz^{k+1})
-\bx_k^\star(\bz^k),
\bx^{k+1}-\bz^k
\right\rangle \\
&\qquad
+\frac{\beta_k(2-\beta_k)}{2s_k}
\|\blmbd^{k+1}-\bmu^k\|^2
+\frac12\left(\frac1{s_k}-\frac1{s_{k+1}}\right)
\|\blmbd^{k+1}-\bmu^{k+1}\|^2.
\end{aligned}
\label{eq:augmented-one-step}
\end{equation}
\end{lemma}

Two gaps remain before $\Phi_k$ admits a genuine descent inequality: the augmented estimate still carries an indefinite cross term and positive multiplier/center-increment terms, and its left-hand side compares $d_k, p_k$ at both endpoints rather than the true $k \to k+1$ potential difference, which requires $d_{k+1}, p_{k+1}$. 

For that, we establish boundedness of the multipliers directly, rather than assuming compactness a priori: both multipliers are controlled by the AL structure and the inner residual, not by projection onto a compact set.
\begin{lemma}
\label{lem:current-multiplier-boundedness}
Suppose Assumptions~\ref{ass:smoothness} and~\ref{ass:h} hold.
Let the iterates be generated by Algorithm~\ref{alg:ddlalm} with
\eqref{eq:basic-parameter-condition} and
\begin{equation}
\label{eq:multiplier-boundedness-condition}
    \tau_k\left(\frac1\rho+\frac1{s_k}\right)\le1.
\end{equation}
Define the constant independent of $k$
\[
    M:=\max\{\|\blmbd^0\|,\|\bmu^0\|,L_h\}.
\]
Then, for every $k\ge0$, the following holds
\[
    \|\blmbd^k\|\le M,
    \qquad
    \|\bmu^k\|\le M.
\]
Moreover, for every $k\ge0$ and every $\bz\in\RR^n$, it holds that
\[
    \|\blmbd_k^\star(\bz)\|\le M.
\]
\end{lemma}

% The augmented one-step estimate also contains the indefinite term
% \[
% \left\langle
% \bx_k(\blmbd^{k+1};\bz^{k+1})-\bx_k^\star(\bz^k),
% \bx^{k+1}-\bz^k
% \right\rangle.
% \]
% Its first factor combines two effects: the dual iterate may differ from
% the dual maximizer, and the saddle-point primal component may move when
% $\bz$ changes. The next lemma quantifies both effects. The first estimate
% relates the dual suboptimality gap $p_k-d_k$ to the corresponding primal
% distance, while the second estimate controls the sensitivity of the
% primal minimizers to the center variable.

\begin{lemma}
\label{lem:saddle-center-stability}
Under the conditions of Proposition~\ref{prop:curvature-smoothness}, for every $k\ge0$, $\blmbd\in\RR^m$, and $\bz\in\RR^n$, it holds that
\begin{equation}
\label{eq:current-saddle-stability}
\|\bx_k(\blmbd;\bz)-\bx_k^\star(\bz)\|^2
\le
\frac{2}{\mu_\gamma}
\bigl[p_k(\bz)-d_k(\blmbd;\bz)\bigr].
\end{equation}
For any $\bz,\bz'\in\RR^n$,
\begin{equation}
\label{eq:current-center-stability}
\max\bigl\{
\|\bx_k(\blmbd;\bz')-\bx_k(\blmbd;\bz)\|,
\|\bx_k^\star(\bz')-\bx_k^\star(\bz)\|
\bigr\}
\le\frac{1}{\gamma\mu_\gamma}\|\bz'-\bz\|.
\end{equation}
\end{lemma}

We now establish an approximate descent inequality for
\Cref{alg:ddlalm} with respect to the Lyapunov function defined in \eqref{eq:lyapunov-decomposition}. We assume 
the parameters are chosen as
\begin{equation}
\label{eq:parameter-choice}
\begin{gathered}
    0<\gamma<\frac{1}{2L_{\calD,\bx}},\qquad
    0<s_k\le s_{k+1},\qquad
    0<\beta_k\le1,\qquad
    0<\eta_k\le\frac{1}{2L_x},\\
    0<\tau_k\le\left(\max\left\{2L_{d,k},\,
        \frac{36G_{\mathcal A}^2}{\eta_k\mu_\gamma^2}\right\}\right)^{-1},\qquad
    0<\alpha_k\le\frac1{36}\min\left\{1,\,
        \frac{\gamma\mu_\gamma\tau_k}{s_k}\right\},
\end{gathered}
\end{equation}
where $L_{d,k}:=\rho^{-1}+s_k^{-1}+G_{\mathcal A}^2/\mu_\gamma$.
Since $L_{d,k}\ge\rho^{-1}+s_k^{-1}$, \eqref{eq:parameter-choice} implies the conditions of Proposition~\ref{prop:curvature-smoothness},
Lemma~\ref{lem:current-multiplier-boundedness},
\eqref{eq:basic-parameter-condition},
\eqref{eq:dual-gain-condition}, and
\eqref{eq:multiplier-boundedness-condition}.

\begin{lemma}
\label{lem:moving-three-value-descent}
\blue{Suppose that Assumptions~\ref{ass:level-boundedness},
\ref{ass:smoothness}, and~\ref{ass:h} hold, and let}
$\{(\bx^k,\blmbd^k,\bz^k,\bmu^k)\}$ be generated by \Cref{alg:ddlalm}
with parameters satisfying~\eqref{eq:parameter-choice}. Then for every
$k\ge0$:
\begin{enumerate}[label=(\alph*)]
    \item for all $\bx\in\RR^n$, $\blmbd\in\RR^m$, and $\bz\in\RR^n$,
    \begin{equation}
    \label{eq:current-lyapunov-lower-bound}
        \Psi_k(\bx,\blmbd;\bz)\ge p_k(\bz)
        \ge\inf_{\bx\in\RR^n}F(\bx)-\frac{(L_h+M)^2}{2s_0};
    \end{equation}
    \item the iterates satisfy
    \begin{equation}
    \label{eq:full-lyapunov-descent}
    \begin{aligned}
        &\Psi_{k+1}(\bx^{k+1},\blmbd^{k+1};\bz^{k+1})
         -\Psi_k(\bx^k,\blmbd^k;\bz^k)\\
        &\qquad\le
         -\frac1{4\eta_k}\|\Delta_{\bx}^{k+1}\|^2
         -\frac1{8\tau_k}\|\Delta_{\blmbd}^{k+1}\|^2
         -\frac{\alpha_k}{2\gamma}\|\bx^{k+1}-\bz^k\|^2
         +\delta_k,
    \end{aligned}
    \end{equation}
    where
    $\delta_k:=10M^2\bigl(s_k^{-1}-s_{k+1}^{-1}\bigr)+20M^2\beta_k/s_k$.
\end{enumerate}
\end{lemma}


% For $s_k=c_s(k+1)^{1/3}$ and
% $\beta_k=\theta_\mu/(k+k_0)$ with $k_0\ge1$, the total error satisfies
% \begin{equation}
% \label{eq:explicit-transport-budget}
% \sum_{k=0}^\infty
% \left[
% 10M^2\left|\frac1{s_{k+1}}-\frac1{s_k}\right|
% +20M^2\frac{\beta_k}{s_k}
% \right]
% \le\frac{(10+80\theta_\mu)M^2}{c_s}.
% \end{equation}

% The augmented estimate adds $2p_k$ to the partial potential in
% \eqref{eq:partial-potential}. It uses the $1/\gamma$-smoothness of $p_k$
% once. The remaining
% inner product is the only new $\bz$-direction error that must be controlled
% by the subsequent stability estimates.

% \begin{lemma}[Saddle-point stability]
% \label{lem:saddle-point-stability}
% Under the conditions of
% Proposition~\ref{prop:curvature-smoothness}, for each $k\ge 0$ and
% each fixed $\bz$, let
% \[
%     \blmbd_k^\star(\bz)
%     :=\argmax_{\blmbd\in\RR^m}d_k(\blmbd;\bz).
% \]
% Then
% \[
%     \bx_k^\star(\bz)
%     =
%     \bx_k\bigl(\blmbd_k^\star(\bz),\bz\bigr),
% \]
% and
% \begin{equation}
% \label{eq:saddle-primal-stability}
% \left\|
%     \bx_k(\blmbd;\bz)-\bx_k^\star(\bz)
% \right\|
% \le
% \frac{G_{\mathcal A}}{\mu_\gamma}
% \left\|
%     \blmbd-\blmbd_k^\star(\bz)
% \right\|.
% \end{equation}
% Moreover,
% \begin{equation}
% \label{eq:dual-gap-to-multiplier-distance}
% \left\|
%     \blmbd-\blmbd_k^\star(\bz)
% \right\|^2
% \le
% 2s_k
% \bigl[
%     p_k(\bz)-d_k(\blmbd;\bz)
% \bigr].
% \end{equation}
% Consequently,
% \begin{equation}
% \label{eq:saddle-point-stability-final}
% \left\|
%     \bx_k(\blmbd;\bz)-\bx_k^\star(\bz)
% \right\|^2
% \le
% \frac{2G_{\mathcal A}^2s_k}{\mu_\gamma^2}
% \bigl[
%     p_k(\bz)-d_k(\blmbd;\bz)
% \bigr].
% \end{equation}
% \end{lemma}

The proof is given in Appendix~\ref{app:curvature-smoothness}.

